\section{From Suspect Propositions to a Doctor of the Church}

\subsection{The condemnations of 1277}

Thomas died before the most famous institutional controversy over the Parisian
Aristotelian debates.  On 7 March 1277---three years to the day after his
death---Bishop Étienne Tempier condemned 219 propositions.  The list does not
name Thomas.  It gathers claims of different origin and coherence; some oppose
positions Thomas expressly rejected, while a significant group touches or was
understood to touch Thomistic theses.  At Oxford later that month, Archbishop
Robert Kilwardby prohibited a smaller set that more directly included
distinctive positions associated with Thomas, especially concerning the unity
of substantial form.

These acts did not amount to one universal condemnation of ``Thomism.''  They
did make Dominican defense necessary and show that Thomas's later authority was
not immediate or uncontested.  In 1325 the bishop of Paris revoked the 1277
condemnation only insofar as it touched or was said to touch Thomas after his
canonization; the act left the propositions to free scholastic discussion and
did not certify each one as true.  A simple fall--vindication plot loses the
textual limits of all three acts.

\subsection{Inquiry and canonization}

Dominican promotion of Thomas's cult and teaching gathered memories, miracle
reports, documents, and petitions.  William of Tocco became a principal
promoter and revised his \emph{Ystoria} during the process.  A papally
authorized inquiry at Naples in 1319 heard witnesses whose proximity varied:
some knew Thomas, some knew companions, and others reported public reputation
or miracles.  A second inquiry at Fossanova in 1321 focused especially on the
tomb and miracles.  The records are precious precisely because they preserve
speaker and chain; extracting every answer into an undifferentiated eyewitness
biography would discard their strongest feature.

John~XXII canonized Thomas at Avignon on 18 July 1323.  Canonization is official
Catholic reception of sanctity and cult.  It is not a modern verdict on every
date in Tocco, a guarantee of every alleged miracle, or approval of each
philosophical conclusion.  The process also demonstrates how memory had become
institutional: a religious order, universities, rulers, witnesses, notaries,
and papal authority all helped make a dead teacher publicly legible as a saint.

The canonization is sometimes adorned with a saying that Thomas had performed
as many miracles as he wrote articles.  The memorable formula entered through
a later printed layer rather than the controlled base witness.  It should not
replace the actual process or make intellectual abundance itself a miracle
dossier.

\subsection{Doctor, patron, and many Thomisms}

Pius~V declared Thomas a Doctor of the Church on 11 April 1567 in
\emph{Mirabilis Deus}.  The title recognizes the eminent service of his
teaching to the Church.  It does not make him a source of revelation or require
that every Catholic school adopt one later commentator's reading.  The Piana
edition of his works and the great commentarial traditions helped create
``Thomism,'' but Dominican, Jesuit, Carmelite, analytic, transcendental, and
other appropriations have not been one method or list of conclusions.

Leo~XIII's \emph{Aeterni Patris} (1879) made Thomas central to a renewal of
Christian philosophy, and the Leonine edition gave that renewal a critical
textual task.  In 1880 \emph{Cum hoc sit} named him patron of Catholic schools.
Seminary curricula, manuals, the modern Thomist revival, and twentieth-century
debate magnified his institutional presence.  Manual abbreviation sometimes
made his thought look more closed than the objections, authorities, and
development of the original works.

Paul~VI's \emph{Lumen Ecclesiae} (1974) commemorated the seventh centenary of
Thomas's death while explicitly allowing the legitimacy of other Catholic
schools.  John Paul~II and Benedict~XVI praised the harmony of faith and reason
and the theological orientation of his wisdom.  In 2025 Leo~XIV described
Thomas and John Henry Newman together as co-patrons of the Church's educational
mission.  The mutable wording matters: Thomas should not be called the sole
current patron on the strength of the 1880 act alone.

His present Roman memorial is 28 January, associated with the medieval
translation of his relics to Toulouse; 7 March remains his death date.  A feast
orders ecclesial remembrance.  It does not settle a birth year, authenticate a
portrait, or turn all later scholastic conflict into disobedience.

\evidenceaxis{The Catholic Church canonized Thomas in 1323 and declared him a
Doctor in 1567.}{The canonization proceedings and bull, \emph{Mirabilis Deus},
and later papal acts are official evidence for their respective acts.}{Their
objects are sanctity, cult, eminent teaching, patronage, and recommendation;
they do not remove textual variants, historical limits, or legitimate Catholic
schools.}
